\documentclass[11pt]{ctexart}
\usepackage[a4paper,margin=1in]{geometry}
\usepackage{longtable,booktabs}
\title{Geant4-Agent 回归测试报告（中文）}
\author{自动化评测流程}
\date{2026-03-07\_122142}
\begin{document}
\maketitle
\section{统一 Pass/Fail 总表}
\begin{longtable}{p{0.06\linewidth}p{0.08\linewidth}p{0.07\linewidth}p{0.08\linewidth}p{0.07\linewidth}p{0.07\linewidth}p{0.07\linewidth}p{0.07\linewidth}p{0.09\linewidth}p{0.06\linewidth}p{0.21\linewidth}}
\toprule
ID & 领域 & 风格 & 参数完整性 & 期望初始 & 实际初始 & 期望最终 & 实际最终 & 缺参(初->终) & 结果 & 失败原因 \\
\midrule
MTM01 & medical & fuzzy & missing & 失败 & 失败 & 通过 & 通过 & 13->0 & 通过 & - \\
MTM02 & aerospace & precise & missing & 失败 & 失败 & 通过 & 通过 & 3->0 & 通过 & - \\
MTM03 & industrial\_ndt & precise & missing & 失败 & 失败 & 通过 & 通过 & 10->0 & 通过 & - \\
MTM04 & physics & precise & missing & 失败 & 失败 & 通过 & 通过 & 8->0 & 通过 & - \\
MTM05 & nuclear\_engineering & fuzzy & missing & 失败 & 失败 & 通过 & 通过 & 13->0 & 通过 & - \\
MTM06 & security\_screening & fuzzy & missing & 失败 & 通过 & 通过 & 通过 & 0->0 & 失败 & initial\_complete\_mismatch, missing\_not\_reduced\_after\_followups \\
MTM07 & semiconductor & precise & missing & 失败 & 失败 & 通过 & 通过 & 5->0 & 通过 & - \\
MTM08 & education & fuzzy & missing & 失败 & 失败 & 通过 & 通过 & 11->0 & 通过 & - \\
MTM09 & environmental\_radiation & fuzzy & missing & 失败 & 失败 & 通过 & 通过 & 13->0 & 通过 & - \\
MTM10 & high\_energy\_physics & precise & missing & 失败 & 失败 & 通过 & 通过 & 10->0 & 通过 & - \\
MTM11 & materials\_update & precise & full & 通过 & 通过 & 通过 & 通过 & 0->0 & 通过 & - \\
MTM12 & overwrite\_reject & precise & full & 通过 & 通过 & 通过 & 通过 & 0->0 & 失败 & expected\_final\_mismatch:materials.selected\_materials.0, expected\_final\_mismatch:output.format \\
\bottomrule
\end{longtable}
\section{基准对比}
\begin{itemize}
\item pass\_rate=0.833
\item initial\_alignment\_rate=0.917
\item final\_alignment\_rate=1.000
\item avg\_turns\_per\_case=2.667
\item missing\_reduction\_rate=1.000
\item confirm\_apply\_success\_rate=1.000
\item rollback\_after\_confirm\_rate=0.000
\item internal\_leak\_turn\_rate=0.094
\item repeated\_assistant\_case\_rate=0.000
\item stale\_confirm\_case\_rate=0.000
\item delta\_pass\_rate=-0.117
\item regressions=无
\item improvements=无
\end{itemize}
\appendix
\section{附录：完整对话与内部处理链（自然语言）}
\subsection{案例 MTM01（medical）}
\textbf{对话：}
\begin{itemize}
\item 用户：???????????????????????????????
\item 系统：您想创建什么形状的几何体呢？如果是盒体的话，它的长度（X方向尺寸）是多少？
\item 用户：? 300 mm x 300 mm x 300 mm ?????? G4\_WATER?source.type=beam?particle=gamma?energy=6 MeV?position=(0,0,-400) mm?direction=+z?physics=QBBC?output=root?
\item 系统：配置已完成。当前已确认的设置有：物理列表选择依据、物理列表选择来源、材料选择依据、材料选择来源。
\end{itemize}
\textbf{内部链路（自然语言）：}
\begin{itemize}
\item 第 1 轮：动作=ask\_clarification，阶段=geometry，闭环状态=False。
\item 缺失字段：geometry.structure、geometry.params.module\_x、geometry.params.module\_y、geometry.params.module\_z、materials.volume\_material\_map、source.type、source.particle、source.energy、source.position、source.direction、physics.physics\_list、output.format、output.path。
\item LLM 使用=True；回退=无；后端=llm\_semantic\_frame+runtime\_semantic。
\item 第 2 轮：动作=finalize，阶段=materials，闭环状态=True。
\item 缺失字段：无。
\item LLM 使用=True；回退=无；后端=llm\_slot\_frame+runtime\_semantic。
\end{itemize}
\subsection{案例 MTM02（aerospace）}
\textbf{对话：}
\begin{itemize}
\item 用户：I need a first-pass proton irradiation setup for a silicon device; keep it simple for now.
\item 系统：Geometry: updated geometry structure, box size along x, box size along y. Materials: updated materials.selected\_materials, volume-to-material binding, material selection source. Source: updated source type, particle type, source energy. Physics: still needs physics list. Output: still needs output format, output path.
\item 用户：Use a silicon box 50 mm x 50 mm x 5 mm, beam source proton 100 MeV from (0,0,-100) mm along +z, physics FTFP\_BERT, output json.
\item 系统：You've already set a geometry size. This change would overwrite it.

Pending change:
- Box size along x: 1000.0 → 50.0
- Box size along y: 1000.0 → 50.0

Reply **'confirm'** to apply this change, or **'keep existing'** to keep the current value.
\item 用户：confirm
\item 系统：Your configuration is complete. The following items have been confirmed: the box size along all three axes (x, y, z) and the output path.
\end{itemize}
\textbf{内部链路（自然语言）：}
\begin{itemize}
\item 第 1 轮：动作=summarize\_progress，阶段=materials，闭环状态=False。
\item 缺失字段：physics.physics\_list、output.format、output.path。
\item LLM 使用=True；回退=无；后端=llm\_slot\_frame+runtime\_semantic。
\item 第 2 轮：动作=confirm\_overwrite，阶段=source，闭环状态=False。
\item 缺失字段：无。
\item LLM 使用=True；回退=无；后端=llm\_slot\_frame+runtime\_semantic。
\item 第 3 轮：动作=finalize，阶段=physics，闭环状态=True。
\item 缺失字段：无。
\item LLM 使用=True；回退=无；后端=llm\_semantic\_frame+runtime\_semantic+confirmed\_pending\_overwrite。
\end{itemize}
\subsection{案例 MTM03（industrial\_ndt）}
\textbf{对话：}
\begin{itemize}
\item 用户：????????????????? gamma ??????????
\item 系统：目前我们已经确定了源的设置：使用伽马粒子，并明确了其方向和选择依据。接下来还需要补充几何结构的具体尺寸、材料绑定，以及源的详细参数（如能量和位置）等信息。
\item 用户：???? 40 mm??? 80 mm??? G4\_Fe?source.type=point?particle=gamma?energy=2 MeV?position=(0,0,-200) mm?direction=+z?physics=Shielding?output=csv?
\item 系统：我注意到这将会覆盖一个已经确认的设置。具体修改如下：

**源类型**：从“isotropic”改为“point”

请确认您的选择：
- 如果要应用这个修改，请直接回复“确认”。
- 如果要保留原来的“isotropic”设置，请直接回复“保持原值”。
\item 用户：确认
\item 系统：配置已完成。当前已确认的设置包括：源类型、输出路径、输出格式和物理列表。
\end{itemize}
\textbf{内部链路（自然语言）：}
\begin{itemize}
\item 第 1 轮：动作=summarize\_progress，阶段=geometry，闭环状态=False。
\item 缺失字段：geometry.structure、geometry.params.module\_x、geometry.params.module\_y、geometry.params.module\_z、materials.volume\_material\_map、source.energy、source.position、physics.physics\_list、output.format、output.path。
\item LLM 使用=True；回退=无；后端=llm\_slot\_frame+runtime\_semantic。
\item 第 2 轮：动作=confirm\_overwrite，阶段=materials，闭环状态=False。
\item 缺失字段：无。
\item LLM 使用=True；回退=无；后端=llm\_slot\_frame+runtime\_semantic。
\item 第 3 轮：动作=finalize，阶段=source，闭环状态=True。
\item 缺失字段：无。
\item LLM 使用=True；回退=无；后端=llm\_semantic\_frame+runtime\_semantic+confirmed\_pending\_overwrite。
\end{itemize}
\subsection{案例 MTM04（physics）}
\textbf{对话：}
\begin{itemize}
\item 用户：Gamma point source 1 MeV at (0,0,-100) toward +z; please build the rest.
\item 系统：Geometry: still needs geometry structure, box size along x. Materials: still needs volume-to-material binding. Source: updated source type, particle type, source energy. Physics: still needs physics list. Output: still needs output format, output path.
\item 用户：Set copper box 1 m x 1 m x 1 m, physics FTFP\_BERT, output hdf5.
\item 系统：Your configuration is complete. The physics and material selections have been confirmed.
\end{itemize}
\textbf{内部链路（自然语言）：}
\begin{itemize}
\item 第 1 轮：动作=summarize\_progress，阶段=geometry，闭环状态=False。
\item 缺失字段：geometry.structure、geometry.params.module\_x、geometry.params.module\_y、geometry.params.module\_z、materials.volume\_material\_map、physics.physics\_list、output.format、output.path。
\item LLM 使用=True；回退=无；后端=llm\_slot\_frame+runtime\_semantic。
\item 第 2 轮：动作=finalize，阶段=materials，闭环状态=True。
\item 缺失字段：无。
\item LLM 使用=True；回退=无；后端=llm\_slot\_frame+runtime\_semantic。
\end{itemize}
\subsection{案例 MTM05（nuclear\_engineering）}
\textbf{对话：}
\begin{itemize}
\item 用户：?????????????????????????
\item 系统：你提到的几何形状具体是什么样的呢？比如是盒体、圆柱，还是其他形状？如果是盒体，它在 x 方向上的尺寸大概是多少？
\item 用户：? 200 mm x 200 mm x 200 mm ?????? G4\_Pb?source.type=point?particle=neutron?energy=14 MeV?position=(0,0,-50) mm?direction=+z?physics=QGSP\_BERT\_HP?output=xml?
\item 系统：配置已完成。当前已确认：物理列表的选择依据和来源，以及材料的选择依据和来源。
\end{itemize}
\textbf{内部链路（自然语言）：}
\begin{itemize}
\item 第 1 轮：动作=ask\_clarification，阶段=geometry，闭环状态=False。
\item 缺失字段：geometry.structure、geometry.params.module\_x、geometry.params.module\_y、geometry.params.module\_z、materials.volume\_material\_map、source.type、source.particle、source.energy、source.position、source.direction、physics.physics\_list、output.format、output.path。
\item LLM 使用=True；回退=无；后端=llm\_semantic\_frame+runtime\_semantic。
\item 第 2 轮：动作=finalize，阶段=materials，闭环状态=True。
\item 缺失字段：无。
\item LLM 使用=True；回退=无；后端=llm\_slot\_frame+runtime\_semantic。
\end{itemize}
\subsection{案例 MTM06（security\_screening）}
\textbf{对话：}
\begin{itemize}
\item 用户：Need a quick cargo-screening toy setup with forward photons; details can follow.
\item 系统：Your configuration is complete. The geometry, materials, source, physics, and output settings are all set.
\item 用户：Use steel cylinder target radius 60 mm half-length 120 mm, beam gamma 3 MeV from (0,0,-500) mm to +z, physics Shielding, output root.
\item 系统：This would overwrite an already confirmed value. Pending change: geometry structure: single\_box -> single\_tubs; materials.selected\_materials: ['G4\_Cu'] -> ['G4\_STAINLESS-STEEL']. Reply 'confirm' to apply it, or say 'keep existing' to keep the current value.
\item 用户：confirm
\item 系统：Configuration complete. All confirmed items are set: volume-to-material binding, geometry root volume, output path, half length, and radius.
\end{itemize}
\textbf{内部链路（自然语言）：}
\begin{itemize}
\item 第 1 轮：动作=finalize，阶段=materials，闭环状态=True。
\item 缺失字段：无。
\item LLM 使用=True；回退=无；后端=llm\_slot\_frame+runtime\_semantic。
\item 第 2 轮：动作=confirm\_overwrite，阶段=source，闭环状态=False。
\item 缺失字段：无。
\item LLM 使用=True；回退=无；后端=llm\_slot\_frame+runtime\_semantic。
\item 第 3 轮：动作=finalize，阶段=physics，闭环状态=True。
\item 缺失字段：无。
\item LLM 使用=True；回退=无；后端=llm\_semantic\_frame+runtime\_semantic+confirmed\_pending\_overwrite。
\end{itemize}
\subsection{案例 MTM07（semiconductor）}
\textbf{对话：}
\begin{itemize}
\item 用户：??????????????????????????
\item 系统：目前我们已经完成了几何结构类型和盒体尺寸的设置，以及源的能量、位置和方向。接下来还需要补充：体积与材料的绑定、粒子类型、物理列表、输出格式和输出路径。
\item 用户：?? 20 mm x 20 mm x 1 mm ?????? G4\_Si?source.type=isotropic?particle=neutron?energy=2 MeV?position=(0,0,0) mm?physics=QGSP\_BERT\_HP?output=json?
\item 系统：我注意到这将会覆盖已经确认的设置。具体修改如下：

- 盒体 x 方向尺寸：1000.0 → 20.0
- 盒体 y 方向尺寸：1000.0 → 20.0

请确认您的选择：
- 如果要应用这些修改，请直接回复“确认”。
- 如果要保留原来的数值，请直接回复“保持原值”。
\item 用户：确认
\item 系统：配置已完成。当前已确认的设置包括：源位置、源能量、源类型以及盒体的 z 方向尺寸。
\end{itemize}
\textbf{内部链路（自然语言）：}
\begin{itemize}
\item 第 1 轮：动作=summarize\_progress，阶段=materials，闭环状态=False。
\item 缺失字段：materials.volume\_material\_map、source.particle、physics.physics\_list、output.format、output.path。
\item LLM 使用=True；回退=无；后端=llm\_slot\_frame+runtime\_semantic。
\item 第 2 轮：动作=confirm\_overwrite，阶段=source，闭环状态=False。
\item 缺失字段：无。
\item LLM 使用=True；回退=无；后端=llm\_slot\_frame+runtime\_semantic。
\item 第 3 轮：动作=finalize，阶段=physics，闭环状态=True。
\item 缺失字段：无。
\item LLM 使用=True；回退=无；后端=llm\_semantic\_frame+runtime\_semantic+confirmed\_pending\_overwrite。
\end{itemize}
\subsection{案例 MTM08（education）}
\textbf{对话：}
\begin{itemize}
\item 用户：For classroom demo, I want gamma attenuation in copper; you can ask only essential missing parts.
\item 系统：Geometry: still needs geometry structure, box size along x. Materials: updated materials.selected\_materials, material selection source, material selection reasons; still needs volume-to-material binding. Source: updated particle type, source selection source, source selection reasons; still needs source type, source energy. Physics: updated physics list, physics.backup\_physics\_list, physics selection reasons. Output: still needs output format, output path.
\item 用户：Okay, copper box 100 mm x 100 mm x 100 mm, point gamma 1 MeV at (0,0,-200) mm, direction +z, physics FTFP\_BERT, output csv.
\item 系统：This would change a setting you've already confirmed.

Pending change:  
**Physics list:** Shielding → FTFP\_BERT

Please reply **"confirm"** to apply this change, or **"keep existing"** to keep the current Shielding physics list.
\item 用户：confirm
\item 系统：Your configuration is complete.

The following items are confirmed: physics list, output path, output format, and source direction.
\end{itemize}
\textbf{内部链路（自然语言）：}
\begin{itemize}
\item 第 1 轮：动作=summarize\_progress，阶段=geometry，闭环状态=False。
\item 缺失字段：geometry.structure、geometry.params.module\_x、geometry.params.module\_y、geometry.params.module\_z、materials.volume\_material\_map、source.type、source.energy、source.position、source.direction、output.format、output.path。
\item LLM 使用=True；回退=无；后端=llm\_slot\_frame+runtime\_semantic。
\item 第 2 轮：动作=confirm\_overwrite，阶段=materials，闭环状态=False。
\item 缺失字段：无。
\item LLM 使用=True；回退=无；后端=llm\_slot\_frame+runtime\_semantic。
\item 第 3 轮：动作=finalize，阶段=source，闭环状态=True。
\item 缺失字段：无。
\item LLM 使用=True；回退=无；后端=llm\_semantic\_frame+runtime\_semantic+confirmed\_pending\_overwrite。
\end{itemize}
\subsection{案例 MTM09（environmental\_radiation）}
\textbf{对话：}
\begin{itemize}
\item 用户：??????????????????????????????
\item 系统：你打算创建什么形状的几何体呢？如果是盒体的话，它在 x 方向的尺寸是多少？
\item 用户：? 500 mm x 500 mm x 500 mm ??????? G4\_AIR?source.type=point?particle=gamma?energy=0.5 MeV?position=(0,0,-300) mm?direction=+z?physics=QBBC?output=root?
\item 系统：配置已完成。当前已确认：物理列表的选择依据和来源、材料的选择依据和来源。
\end{itemize}
\textbf{内部链路（自然语言）：}
\begin{itemize}
\item 第 1 轮：动作=ask\_clarification，阶段=geometry，闭环状态=False。
\item 缺失字段：geometry.structure、geometry.params.module\_x、geometry.params.module\_y、geometry.params.module\_z、materials.volume\_material\_map、source.type、source.particle、source.energy、source.position、source.direction、physics.physics\_list、output.format、output.path。
\item LLM 使用=True；回退=无；后端=llm\_semantic\_frame+runtime\_semantic。
\item 第 2 轮：动作=finalize，阶段=materials，闭环状态=True。
\item 缺失字段：无。
\item LLM 使用=True；回退=无；后端=llm\_slot\_frame+runtime\_semantic。
\end{itemize}
\subsection{案例 MTM10（high\_energy\_physics）}
\textbf{对话：}
\begin{itemize}
\item 用户：Set geometry as a sphere target; I have not decided radius or source details yet.
\item 系统：You've set the geometry to a sphere target and defined the root volume. We still need the sphere's radius.

We also need to assign a material to that volume, and define the particle source (its type, particle, and energy). Finally, we'll need to select a physics list and specify the output format and file path.
\item 用户：Sphere target radius 80 mm, lead material, electron point source 5 MeV at (0,0,-200) mm toward +z, physics FTFP\_BERT, output hdf5.
\item 系统：Your configuration is complete. The physics and material selections have been confirmed.
\end{itemize}
\textbf{内部链路（自然语言）：}
\begin{itemize}
\item 第 1 轮：动作=summarize\_progress，阶段=geometry，闭环状态=False。
\item 缺失字段：geometry.params.child\_rmax、materials.volume\_material\_map、source.type、source.particle、source.energy、source.position、source.direction、physics.physics\_list、output.format、output.path。
\item LLM 使用=True；回退=无；后端=llm\_slot\_frame+runtime\_semantic。
\item 第 2 轮：动作=finalize，阶段=materials，闭环状态=True。
\item 缺失字段：无。
\item LLM 使用=True；回退=无；后端=llm\_slot\_frame+runtime\_semantic。
\end{itemize}
\subsection{案例 MTM11（materials\_update）}
\textbf{对话：}
\begin{itemize}
\item 用户：Create a complete setup: 100 mm x 100 mm x 100 mm copper box, point gamma 1 MeV at (0,0,-200) mm toward +z, physics QBBC, output json.
\item 系统：Current committed config is: physics list FTFP\_BERT, output format json, source type point, particle gamma. If you want changes, please specify which field to update.
\item 用户：Change the target material to G4\_Al and change the output format to root.
\item 系统：You're about to overwrite two confirmed settings:
- The target material will change from copper (G4\_Cu) to aluminum (G4\_Al).
- The output format will change from JSON to ROOT.

Please reply **"confirm"** to apply these changes, or **"keep existing"** to keep the current values.
\item 用户：confirm
\item 系统：Configuration complete. The following items are confirmed: volume-to-material binding, output path, output format, and selected materials.
\end{itemize}
\textbf{内部链路（自然语言）：}
\begin{itemize}
\item 第 1 轮：动作=finalize，阶段=materials，闭环状态=True。
\item 缺失字段：无。
\item LLM 使用=True；回退=无；后端=llm\_slot\_frame+runtime\_semantic。
\item 第 2 轮：动作=confirm\_overwrite，阶段=source，闭环状态=False。
\item 缺失字段：无。
\item LLM 使用=True；回退=无；后端=llm\_slot\_frame+runtime\_semantic。
\item 第 3 轮：动作=finalize，阶段=physics，闭环状态=True。
\item 缺失字段：无。
\item LLM 使用=True；回退=无；后端=llm\_semantic\_frame+runtime\_semantic+confirmed\_pending\_overwrite。
\end{itemize}
\subsection{案例 MTM12（overwrite\_reject）}
\textbf{对话：}
\begin{itemize}
\item 用户：????????????100 mm x 100 mm x 100 mm ?????? G4\_Cu??? gamma 1 MeV??? (0,0,-200) mm??? +z?physics=QBBC?output=json?
\item 系统：配置已完成。当前已确认：物理列表选择依据、物理列表选择来源、材料选择依据、材料选择来源。
\item 用户：????? G4\_Al??????? root?
\item 系统：我注意到这次修改会覆盖已经确认的配置项。具体来说：

- **材料**：从 `G4\_Cu` 改为 `G4\_Al`
- **输出格式**：从 `json` 改为 `root`

如果你确定要应用这些修改，请直接回复“确认”；如果想保留原来的值，请回复“保持原值”。
\item 用户：????
\item 系统：我注意到这次修改会覆盖已经确认的配置项：

- **材料**：从 `G4\_Cu` 改为 `G4\_Al`
- **输出格式**：从 `json` 改为 `root`

如果你确定要应用这些修改，请回复“确认”；如果想保留原来的值，请回复“保持原值”。
\item 用户：确认
\item 系统：好的，配置已完成。当前已确认的设置包括：体积与材料绑定、输出路径、输出格式以及所选材料。
\end{itemize}
\textbf{内部链路（自然语言）：}
\begin{itemize}
\item 第 1 轮：动作=finalize，阶段=materials，闭环状态=True。
\item 缺失字段：无。
\item LLM 使用=True；回退=无；后端=llm\_slot\_frame+runtime\_semantic。
\item 第 2 轮：动作=confirm\_overwrite，阶段=source，闭环状态=False。
\item 缺失字段：无。
\item LLM 使用=True；回退=无；后端=llm\_slot\_frame+runtime\_semantic。
\item 第 3 轮：动作=confirm\_overwrite，阶段=physics，闭环状态=False。
\item 缺失字段：无。
\item LLM 使用=True；回退=无；后端=llm\_slot\_frame+runtime\_semantic。
\item 第 4 轮：动作=finalize，阶段=output，闭环状态=True。
\item 缺失字段：无。
\item LLM 使用=True；回退=无；后端=llm\_semantic\_frame+runtime\_semantic+confirmed\_pending\_overwrite。
\end{itemize}
\end{document}